\section{Introduction}
\label{sec:intro}

Multi-task serving and model merging have risen as essential paradigms for deploying specialized deep learning models under strict resource constraints. Rather than maintaining separate base models for multiple disparate tasks, contemporary approaches fine-tune specialized Parameter-Efficient Fine-Tuning (PEFT) adapters, such as Low-Rank Adaption (LoRA) \cite{hu2021lora}, and fuse or route between them at runtime. Dynamic activation-space ensembling \cite{sable2025samplewise, chemmerge2025kinetics} represents a powerful family of these methods, in which intermediate representation vectors are mapped dynamically to ensembling or gating coefficients at each layer of a deep neural network, routing the input token or sample through specialized expert adapters.

\begin{figure*}[t]
  \centering
  \begin{subfigure}[b]{0.48\textwidth}
    \centering
    \includegraphics[width=\textwidth]{results/fig1.png}
    \caption{Performance comparison.}
    \label{fig:performance_comparison}
  \end{subfigure}
  \hfill
  \begin{subfigure}[b]{0.48\textwidth}
    \centering
    \includegraphics[width=\textwidth]{results/layer_trajectory.png}
    \caption{Layer-wise routing coefficient trajectory.}
    \label{fig:trajectory_comparison}
  \end{subfigure}
  \caption{(a) Main empirical performance results across all serving baselines over 10 independent random seeds under orthogonal and overlapping sandbox settings. Capsize bars show standard deviation. (b) Layer-wise coefficient trajectory of a representative sample under routing. Unregularized parametric routing exhibits severe high-frequency oscillations across depth, whereas the proposed CR-Router stabilizes smoothly to a unique fixed-point trajectory.}
  \label{fig:main_figures}
\end{figure*}

Despite its empirical promise, sequential dynamic ensembling introduces a critical, yet frequently overlooked challenge: \emph{sequential routing jitter}. Because the routing decision at layer $l$ is a function of the intermediate feature $h^{(l-1)}$---which itself depends on the routing decisions and adapter blend at previous layers---this iterative feedback loop forms a highly complex, discrete-time dynamical system across the depth of the network. When left unregularized, the layer-wise routing coefficients exhibit high-frequency oscillations from layer to layer, causing the representation trajectory to undergo violent functional shifts. This destabilization degrades the model's joint classification capability and makes the dynamic router highly susceptible to transductive overfitting, especially in data-scarce settings where calibration data is severely limited (e.g., only 16 samples per task).

To mitigate this instability, prior works such as \cite{chemmerge2025kinetics} introduce heuristics inspired by continuous chemical kinetics to penalize temporal differences in gating decisions. However, these methods are primarily empirical and lack a solid mathematical foundation to explain \emph{why} representation flow converges stably across depth. Furthermore, they do not resolve the severe overfitting behavior that unregularized routers exhibit under sample scarcity. We argue that empirical ensembling success is fragile without formal convergence guarantees and rigorous theoretical parameter bounds.

In this work, we approach this problem from a rigorous mathematical perspective. We model the feedforward propagation of representations and routing coefficients across sequential network layers as an iterative mapping on a Banach space. Using this framework, we analyze the functional properties of this feedback operator and derive a novel Lipschitz bound on the joint layer-wise representation-routing map. Under Banach's Fixed-Point Theorem, we show that we can constrain this joint map to be a strict contraction by penalizing the spectral norm of the routing projection matrices and bounding the inverse routing temperature parameters. 

Based on this theoretical insight, we introduce the \textbf{Contraction-Regularized Router (CR-Router)}. CR-Router implements a theoretically sound regularized objective function containing both Frobenius norm penalties (as analytical upper bounds on the spectral norm) on the routing heads and inverse routing temperature penalties. 
While this framework mathematically guarantees strict contraction and convergence for systems with scaled residuals, direct application to frozen pre-trained residual models (such as deep Transformers) requires a theoretical relaxation. Specifically, we formalize \textbf{Update-Space Quasi-Contraction}, where we trade absolute representational convergence to a global Banach fixed point in order to preserve the frozen backbone's pre-trained capacity. We show that by bounding the update-space Lipschitz constant ($L_{U_l} < \epsilon$), we successfully stabilize the gating trajectories and prevent routing jitter without modifying the frozen model parameters.

We evaluate CR-Router inside a 14-layer, 192-dimensional Analytical Coordinate Sandbox (ICS) \cite{sable2025samplewise} containing specialized adapters for MNIST, Fashion-MNIST, CIFAR-10, and SVHN. Across 10 independent random seeds under extreme data scarcity, CR-Router completely eliminates layer-wise routing jitter and recovers a stellar classification accuracy of \textbf{53.35\% $\pm$ 3.84\%} in perfectly orthogonal subspaces (Experiment 1), outperforming the unregularized linear router by \textbf{18.62\%} absolute. Crucially, in overlapping task subspaces where adjacent tasks share 48 dimensions of overlap (Experiment 2), static Uniform Merging collapses to a near-random 27.48\% classification accuracy due to representation cross-talk, whereas CR-Router recovers a strong \textbf{43.48\% $\pm$ 4.70\%} classification accuracy. This represents a massive absolute improvement of \textbf{16.00\%} over static Uniform Merging and \textbf{12.86\%} over unregularized routing, demonstrating the absolute necessity of dynamic ensembling and contraction regularizers under representational overlap.

To address standard seed sensitivity under representational overlap, we introduce \textbf{Centroid-Based Routing Warm-Starting} as an elegant initialization strategy. By warm-starting the linear routing parameters using task-specific centroids extracted from the extremely scarce calibration split, we provide a strong geometric prior from the start of optimization. This guides early gradient steps directly into stable, task-aligned attraction basins, significantly reducing standard deviation and stabilizing convergence while preserving strict contraction bounds. We recommend this heuristic as the primary initialization strategy for real-world deployments under sample scarcity.

Our core contributions are:
\begin{enumerate}
  \item \textbf{Dynamical Systems Formulation:} We formalize the feedforward feature-coefficient propagation of sequential ensembling as a discrete-time dynamical system.
  \item \textbf{Theoretical Generalization and Lipschitz Bounds:} We prove a novel, strict bound on the Lipschitz constant $L_{T_l}$ of the joint layer-wise mapping.
  \item \textbf{CR-Router Design:} We design a theoretically grounded ensembling router with explicit spectral and inverse temperature penalties to enforce contraction properties ($L_{T_l} < 1$).
  \item \textbf{High-Resolution Sensitivity Analysis:} We map the stability-accuracy trade-off curve with a fine-grained grid sweep over the joint penalty $\lambda$, validating our three proposed online, label-free tuning heuristics (Gating Depth-Variance, Shannon Gating Entropy, and Running Gating Lipschitz Bound).
  \item \textbf{Adaptive Test-Time Annealing Breakthrough:} We propose and validate \emph{Adaptive Test-Time Temperature Annealing} as an elegant post-hoc solution to the "expert dilution" dilemma, successfully decoupling training-time optimization stability from test-time representation sharpness and unlocking massive performance gains (+8.00\% absolute classification accuracy).
  \item \textbf{Centroid-Based Routing Warm-Starting:} We propose and evaluate an elegant initialization heuristic that warm-starts the routing weights using task-specific centroids of the calibration split. We highlight this as the primary recommended initialization method for practitioners deploying in extreme data-scarce settings, showing that it effectively mitigates seed variance and guides optimization into stable, task-aligned basins.
\end{enumerate}
